% Separate post-CP19 candidate. Background data only; no representation count.
\subsection{Moving line backgrounds without matrix-valued basis variables}

\begin{lemma}[Candidate: compressed minors give a bounded rational chart cover]
\label{D:compressed-minor-background-charts}
Let $X$ be a finite-type scheme over an algebraically closed field
$k$, with $\dim X\leq b$, and let $A$ be an $M\times N$ matrix
of regular functions on $X$. On its scheme-theoretic rank-$\rho$
locus, there are $b+1$ pairs of constant matrices
\[
 U_\alpha\in\operatorname{Mat}_{\rho\times M}(k),\qquad
 V_\alpha\in\operatorname{Mat}_{N\times\rho}(k)
\]
such that the opens $\det(U_\alpha A V_\alpha)\ne0$ cover.
On each open, image, kernel, and cokernel have explicit rational
bases, with no matrix-valued moduli variables. If the entries
of $A$ are rational functions in $b$ ambient coordinates, of
numerator and denominator degree at most $E$, these bases have
degrees bounded by a polynomial in $M,N,E$ after clearing
the existing denominators.

Here the rank-$\rho$ locus means that all $(\rho+1)$-minors
vanish and at least one $\rho$-minor is invertible. These are
scheme conditions, including over nonreduced coefficient rings.
For $\rho=0$ no compression is needed.
\end{lemma}

\begin{proof}
On the reduced rank-$\rho$ locus fix a geometric point $x$.
The polynomial $\det(UA(x)V)$ in the entries of $U,V$ is not
identically zero. Its zero set is therefore a proper hypersurface
in the affine space of compression pairs. For $b+1$ independent
pairs, the incidence set on which all these determinants vanish
has dimension at most
\[
 b+(b+1)(\dim\operatorname{Mat}_{\rho\times M}
          +\dim\operatorname{Mat}_{N\times\rho}-1),
\]
strictly smaller than the dimension of the space of all pairs.
Its constructible image has closure of smaller dimension.
Choose constants outside that closure. The resulting opens
cover all geometric points and hence the entire rank locus,
including its nilpotent structure.

Write $U,V$ for one selected pair and $B=UAV$. The constant
matrices $U,V$ have full rank. Choose a constant complement
$W$ to the columns of $V$, so that $[V\ W]$ is invertible.
Then $AV$ is an image basis and
\[
                       W-VB^{-1}UAW                      \tag{B.MB.1}
\]
is a kernel basis. Indeed, the rank equations and the invertible
minor give $A=AVB^{-1}UA$ over every coefficient ring on this
chart. Likewise the fixed space $\ker U$ is a complement to
$\operatorname{im}A$, giving a cokernel basis and a projection.
The inverse $B^{-1}=\operatorname{adj}(B)/\det B$ proves the
degree assertion. This is a finite polynomial calculation;
no choice of a basis over a residue field is being extended
without its determinant being inverted.
\end{proof}

\begin{proposition}[Candidate: rational moving-line cohomology and Cech backgrounds]
\label{D:moving-line-polynomial-background}
Fix a constant $c$ and an integer $h_0$. Let $m,s\geq1$, and
let $X$ be a smooth projective connected curve over an
algebraically closed field $k$, of genus at most $cm$, with a
fixed finite separable map $x:X\to\mathbb P^1$ of degree at
most $cm$. Fix at most $h_0$ moving line bundles with prescribed
degrees of absolute value at most $cms$.

There are ordered-divisor parameter charts with $O(ms)$ ambient
coordinates, equations of degree polynomial in $m,s$, and open
conditions, having the following properties.
\begin{enumerate}
\item Every collection of the prescribed lines over $k$ occurs.
For every finite local Artinian $k$-algebra with residue field $k$,
each collection of line-bundle deformations lifts to such a chart
after a choice of a geometric divisor presentation. The chart
formulas themselves are valid over arbitrary coefficient rings.
\item Consider $O(m)$ line families obtained from this fixed
number of moving lines by tensor products of bounded length,
bounded positive or negative powers, one-step Frobenius in a
fixed characteristic, and fixed deck translations preserving
$x$, together with fixed line twists of degree $O(ms)$.
Their section and cohomology complexes, evaluations, and fixed
finite-pole restriction maps are represented by matrices of
sizes and rational-entry degrees polynomial in $m,s$.
No entries of these matrices are additional moduli coordinates.
\item The base admits a cover by locally closed rational charts
on which the relevant cohomology modules are free and all the
specified matrices have rational splittings. The number of
choices is at most
\[
                      \exp\bigl(O(m\log(ms+1))\bigr).     \tag{B.MB.2}
\]
The equations of these charts have degree polynomial in $m,s$.
They can include all minors required for exact-rank conditions;
no bound on the number of those minors is asserted or needed.
\item On these charts one can choose actual representative,
projection, and homotopy maps for the two-affine line Cech
complexes, compatibly through all higher pole ranges. Through
any prescribed pole range polynomial in $m,s$, all their
coefficients are rational functions of degree polynomial in
$m,s$. This construction introduces no new charts at each
successive pole order.
\end{enumerate}
All implicit polynomials and constants depend only on $c,h_0$,
the bounded tensor-word pattern, and the specified polynomial
pole cutoff. Fixed curve coefficients and fixed line twists are
background constants. For a chain of $f$ phases with a bounded
number of moving lines per phase and only bounded-range
interactions, the ambient background coordinates number
$O(fms)$ and the chart count is
$\exp(O(fm\log(ms+1)))$.

The rank charts cover geometric points. An arbitrary noncoflat
Artinian family is not asserted to factor through one exact-rank
chart. All-ring splitting statements apply to points of the
stated chart schemes. This distinction does not affect the
Artinian surjectivity of the preceding ordered-divisor charts.
\end{proposition}

\begin{proof}
Put $t=\deg x$ and $H=x^{-1}(\infty)$; choose the coordinate
so that $H'=x^{-1}(0)$ is disjoint from $H$. These divisors
need not be reduced. Both have degree $t$.

\emph{A plane chart of controlled degree.}
Choose the least positive $a$ with $at\geq2g(X)+1$. The complete linear series
of $aH$ is very ample. A general $z\in H^0(X,\mathcal O(aH))$
is primitive for the separable extension $k(X)/k(x)$.
Indeed, this linear series generates $k(X)$, and the finitely
many proper intermediate fields of a finite separable extension
cannot contain a general member of its $k$-linear section space.
We may also require the allowed pole order at every point of
$H$. Its minimal polynomial is
\[
 F(x,z)=z^t+c_1(x)z^{t-1}+\cdots+c_t(x),
               \qquad c_i(x)\in k[x],\quad\deg c_i\leq ai.
                                                               \tag{B.MB.3}
\]
Integrality away from infinity gives polynomial coefficients;
their valuations at infinity give the displayed degree bounds.
Since $at\leq2g(X)+1+t$, both degrees of $F$ are $O(m)$.
Its discriminant $\Delta(x)$ is nonzero and has degree $O(m^2)$.
After removing its zeroes, $x=0$, and finitely many other fixed
fibres when required below, the normalization agrees with
\[
             k[x,z,1/(x\Delta)]/(F).                    \tag{B.MB.4}
\]
In fact this algebra is finite etale over the indicated open
of the $x$-line, hence normal, and has the same function field
as $X$. Denote the resulting open curve by $X^\circ$.

\emph{Ordered divisors, including Artinian lifting.}
Choose $E_0=eH$ so that, for every prescribed degree $\ell$,
\[
          d=\ell+\deg E_0\geq2g(X)+1,
                  \qquad d=O(ms).                       \tag{B.MB.5}
\]
Every line $L$ of degree $\ell$ has a presentation
$L=\mathcal O_X(D-E_0)$ with
$D=P_1+\cdots+P_d$, where the $P_i$ belong to $X^\circ$
and their $x$-coordinates are pairwise distinct. To verify
this uniformly, use the very ample line $L(E_0)$. Avoiding
the finite omitted set is open on its section space. A double
zero, or two zeroes at different points with the same
$x$-coordinate, imposes two independent linear conditions
while those points range over a one-dimensional family.
The corresponding bad incidences have smaller dimension.
A general section therefore has the asserted zero divisor.

Each $P_i$ is given by two coordinates $(a_i,b_i)$ satisfying
$F(a_i,b_i)=0$. The conditions
$\Delta(a_i)a_i\ne0$ and $a_i\ne a_j$ are open. This is an
$O(ms)$-coordinate parameter space. For a finite local Artinian
base $T$, vanishing of $H^1(L(E_0))$ on the closed fibre lifts
a chosen section successively through all small extensions.
The lifted zero divisor is finite etale over $T$: at each
closed zero the relative derivative is a unit. Since the
residue field is algebraically closed, it splits uniquely
into sections after the closed zeroes have been ordered.
They retain all the open conditions. The regular section
identifies $L_T(E_0)$ with $\mathcal O(D_T)$. This proves the
stated Artinian lifting assertion directly, without requiring
an unproved surjectivity for cohomology-rank strata.

\emph{Polynomial linear equations for sections.}
Let $h_D(x)=\prod_i(x-a_i)$. Its divisor is
$D+D^{\mathrm c}-dH$, where $D^{\mathrm c}$ is the complement
of the selected points in the corresponding unramified fibres.
Multiplication by $h_D$ identifies sections of $L(qH)$ with
the subspace of
\[
 H^0\bigl(X,\mathcal O((d+q)H-E_0)\bigr)                 \tag{B.MB.6}
\]
that vanishes on $D^{\mathrm c}$. The ambient space is fixed;
for all pole ranges used here its dimension is polynomial
in $m,s$, typically $O(m^2s)$. In the finite etale fibre algebra
$T[z]/F(a_i,z)$, complementary-fibre vanishing is exactly
\[
                  (z-b_i)v(a_i,z)=0.                     \tag{B.MB.7}
\]
Indeed the simple root $b_i$ defines a direct factor $T$,
and the annihilator of $z-b_i$ is precisely that factor.
Thus (B.MB.7) is exact over nonreduced coefficient rings too.
Negative moving points impose evaluation-zero conditions.
For bounded multiplicities, use the corresponding finite
jets: the simple-root factor and its idempotent lift uniquely
modulo $(x-a_i)^c$. Successive coefficient comparison divides
only by the unit $F_z(a_i,b_i)$, and gives rational formulas
of polynomial degree for each bounded $c$.

We give a degree check for the evaluations, rather than assume
that arbitrary regular functions have low-complexity expressions.
The trace-dual basis to $1,z,\ldots,z^{t-1}$ has the form
$Q_j(z)/F_z(z)$, with $\deg_zQ_j<t$ and coefficients drawn
from those of $F$. If a function $v$ has poles only at $H$,
of order at most $qH$, then
\[
 v=\frac{\sum_{j=0}^{t-1}v_j(x)z^j}{F_z(x,z)},
                  \deg v_j\leq q+O(m).                  \tag{B.MB.8}
\]
For clarity, $F_zv$ belongs to $k[x,z]$: trace integrality
puts every element of the normalization in the trace dual
$F_z^{-1}k[x,z]$. The coefficient formula is
$v_j=\operatorname{Tr}(vQ_j)$, and the pole bound gives its
degree. This remains valid with inseparable residue at branch
fibres; only separability of $k(X)/k(x)$ is used.
Fixed finite poles can first be cleared by a fixed polynomial
in $x$. The total degree needed is polynomial in the degree
of their fixed divisor. Evaluating (B.MB.8), or reducing it
modulo $F(a_i,z)$, uses only inverses of $F_z$ or its resultant,
which are units on the chosen open. Adjugate formulas give
polynomial-size matrices of polynomial rational degree.

The same construction applies to any bounded tensor word in
the moving divisors and fixed twists. Representatives of the
fixed twists may be chosen with positive and negative divisor
degrees $O(ms+m)$ by Riemann--Roch. Choose all moving points
away from their finitely many fibres. Multiplication by the
bounded powers of the $h_D$ clears the moving poles; selected
and complementary fibre jets impose all remaining conditions.
One-step Frobenius replaces a moving divisor by $pD$ and raises
coefficient coordinates to their $p$th powers, with $p$ fixed.
A deck translation preserves $x$; its effect on $z$ has the
same pole bound as $z$ and hence a representation (B.MB.8)
of polynomial degree. Thus the listed operations retain the
claimed bounds.

To compute an actual cohomology complex, choose a sufficiently
positive fixed pole divisor so that the ambient line has
vanishing $H^1$, and use its exact section sequence with the
finite divisor of imposed zeroes. This is a two-term complex
of finite free modules before taking cohomology. Alternatively
one may use the finite Cech sequence below. All evaluation
and principal-part maps are actual sheaf maps, so their
matrix models commute with arbitrary coefficient change.
Riemann--Roch and duality bound the total cohomology of each
line of degree $O(ms)$ by $O(ms)$
\citep[Tags 0BS2 and 0BS6]{StacksBNegativeCurve}.

\emph{Rank charts and their number.}
Every raw matrix just constructed has dimensions and entry
degrees bounded by one polynomial in $m,s$. Choose its rank,
impose all next-size minors, and apply
Lemma~\ref{D:compressed-minor-background-charts} to cover
that exact-rank locus by at most $b+1$ rational-splitting
charts, where $b=O(ms)$ is the ambient parameter dimension.
A bounded number of such reductions handles a single
two-term complex and its prescribed evaluation maps.
Each line is processed independently from its raw evaluation
matrices; there is no chain of $O(m)$ successive matrix
inversions feeding into later matrices.
Thus every resulting rational formula still has polynomial
degree. There are polynomially many choices of rank and
$O(ms)$ compression charts for each of $O(m)$ matrices,
which proves (B.MB.2). All defining minors may be retained:
their degrees are polynomial even though their number can
be large. Formula (B.MB.1) shows that the splittings are exact
over every ring satisfying the chart equations.

\emph{Compatible finite-pole Cech homotopies.}
Use $X=(X\setminus H)\cup(X\setminus H')$. For one of the
line families $M$, choose $b_0$ so that
\[
 H^1\bigl(M((b_0-1)H)\bigr)
   =H^1\bigl(M((b_0-1)H')\bigr)=0                       \tag{B.MB.9}
\]
on every fibre. It suffices that both degrees exceed
$2g(X)-2$, so $b_0$ is polynomial in $m,s$.
For $b\geq b_0$, the finite section complex
\[
 C_b^0=H^0(M(bH))\oplus H^0(M(bH')),
       \qquad C_b^1=H^0(M(bH+bH'))                       \tag{B.MB.10}
\]
computes $R\Gamma(M)$. This follows from the exact sheaf
sequence
$0\to M\to M(bH)\oplus M(bH')\to M(bH+bH')\to0$,
since $H,H'$ are disjoint and the two first cohomology groups
vanish. The high-degree section modules and evaluations
are free with base change on the chosen rational charts.

Split $C_{b_0}$ once, by the preceding matrix construction,
onto its free cohomology. Also choose once lifts
$\alpha_1,\ldots,\alpha_t\in H^0(M(b_0H))$ of a basis of
the principal parts at the last $H$-layer, and analogous
$\beta_1,\ldots,\beta_t\in H^0(M(b_0H'))$ at $H'$.
Their existence and coefficient-base change follow from
(B.MB.9). The relevant line on the fixed fibres is independent
of the moving divisor, which is disjoint from those fibres;
one may fix its $k$-basis in advance.

For $b>b_0$, the new principal-part layers are lifted by
\[
          x^{b-b_0}\alpha_i,
                         \qquad x^{-(b-b_0)}\beta_i.     \tag{B.MB.11}
\]
The first family has poles only at $H$, and the second only
at $H'$. In the quotient $C_b/C_{b-1}$ these give matching
bases in degrees zero and one, with the Cech differential
an isomorphism. Hence the complex is the direct sum of
$C_{b_0}$ and compatible contractible pairs, one set for
each new pole order. Define the homotopy on each pair by
that inverse differential, and set its projection to
cohomology equal to zero. The result is a strong deformation
retract of the full two-affine Cech complex, compatible with
every coefficient change on the chart. It satisfies
\[
 \rho\iota=1,\quad \iota\rho=1-dh-hd,\quad
                 h^2=h\iota=\rho h=0.                   \tag{B.MB.12}
\]
It does not increase the pole range above
$\max(b,b_0)$ for an input in $C_b$.

The formulas through a given pole bound are effective rational
formulas in the original point coordinates. Expand at the
two fixed fibres and eliminate the highest principal part,
using (B.MB.11), then repeat. The initial lift matrices are
fixed rational matrices on the chart. Each elimination is
linear in the current coefficients and uses only these fixed
matrices. Thus the degrees grow at most linearly in the
number of pole steps, times their initial degree. The jets
of $h_D^{\pm1}$ have polynomial degree through any polynomial
cutoff: at infinity use
$h_D/x^d=\prod_i(1-a_i/x)$, and at zero use the invertible
constant $h_D(0)$. The same statement holds for bounded powers
and the fixed twists. Consequently products of polynomially
bounded expanded word length, with a bounded number of one-step
Frobenius operations, have polynomial rational coefficient degree
through a polynomial pole cutoff. In the normal-layer application
the word length is at most $am+O(1)$ and there is one fixed-$p$
pullback. This does not assert the same bound for arbitrary
repeated-squaring circuits.
This proves the final assertion, including actual homotopy
data rather than only dimensions or residue-field bases.

The Cech computation and its coefficient description use
the exact affine-intersection and base-change statements in
\citep[Tags 0BDX, 01XD, 02KE]{C_StacksFiniteCechHom}; the
rational matrices, compression argument, and pole-compatible
splittings above are constructed explicitly.
\end{proof}

\begin{remark}[Reference lifts and the limits of the background assertion]
In the one-curve formal surface setting with finite etale
spectral covers, the ordered moving divisor also has an actual
formal lift. Shrink one affine open of the original smooth
curve so that its normal line is trivial. Formal smoothness
and the affine square-zero lifting criterion give an
isomorphism of its formal neighbourhood with the product
with a formal disc. A finite etale spectral cover pulls back
to the corresponding product. Choose all moving divisor
points in this open, as permitted in the proof, and take their
constant transverse Cartier sections. Their divisor lines
glue with the trivial line away from their support; subtract
a fixed formal lift of $E_0$. This produces an actual formal
reference line with the same background point parameters.

On a retained finite normal order, its moving transition
factors are the products of the local equations of these
sections and their inverses. Taylor expansion and inversion
use the same unit denominators as above; through order $am$
their degree in the moving point coordinates is polynomial
in $m,s$. Fixed formal-neighbourhood coefficients remain
background constants. This observation does not assert a
normal-pole bound for arbitrary additional anchor data:
the finite cochain ranges needed for those data must still
be established in the formal object construction.

The proposition provides rational backgrounds and coflat
rank charts. It does not construct the full graded formal
bundle atlas, prove rigidity of its unit objects, or bound
representations. In particular no rank-stratum statement
here repairs the failed open-fullness inference recorded
in the separate genus-two research note.
\end{remark}
